\chapter{Introduction}

\gls{cnn} perform well at detecting objects in controlled environments. But moving to the real world, high-resolution processing is still a challenge when processing on edge system. Most small processors can’t store and process large images directly. A common way to handle this is to shrink the image before processing. This causes the loss of smaller details that are valuable for accurate detection of objects.This project is heavily influenced from field of precision agriculture. Where, large areas are analysed by breaking them into smaller sections and processing each one separately.Here, We applied the same idea, combined with distributed edge computing system inorder to build a prominent solution for real-time processing \gls{hpd}. We chose \gls{hpd} as our focus to solve a specefic problem as Detecting people in large, open spaces requires fine visual details. Traditional downsampling destroys that granularity.
Our approach splits a high resolution image into smaller square tiles then each tile will be  processed by a microcontroller spreading the workload across many similar devices instead of relying on a single powerful one. As a result, we preserve the image quality needed to detect people reliably while processing on alower power devices.
This method also has practical benefits beyond accuracy as  all processing happens locally on the device without needing the internet or the cloud compute. It reduces latency and preserves privacy.


\section{Background}

% Test abbreviations and symbols
Common object detection models like You Only Look Once (YOLO) and Faster architecture like R-CNN are designed to work with small input images with sizes like  416×416 or 640×640 pixels. However, cameras used in surveillance or agricultural automation can capture images as large as 4000×3000 pixels which is nearly ten times larger than what these models expect. 


\subsection*{The Problem of Vanishing Detail}
Surveillance systems often capture very large images but most CNN models only accept small, fixed size inputs. To bridge this gap, the, image is shrunk down before being fed into the model which is known as the Spatial Information Attrition problem which causes serious data loss. 
Consider a frame that is 4000 pixels wide being reduced to 400 pixels. That’s a 90\% reduction in width and 99\% of the total pixel data is thrown away.

Small objects suffer the most from this process. Take a person standing at a distance. They might occupy a 40×80 pixel region in the original image but after downsampling, that same person is reduced to just 4×8 pixels which is roughly the size of a thumbnail on a postage stamp

At that size, a \gls{cnn} cannot do its job properly. Every CNN has a minimum region it needs to see in order to pick up useful features which is called the Effective Receptive Field (ERF) \cite{luo2016understanding}. A 4×8 pixel blob falls well below that minimum which does not have enough details to identify limbs, a torso, or any other feature that distinguishes a person from background noise. The model ends up trying to classify a blurry smudge, which can look suprisingly identical to a shadow or a bush.

\subsection*{Lessons from High-Precision Agriculture}
In precision agriculture, satellites and drones are used to capture very high resolution image of fields. Researchers then analyse these images to spot individual plants,plots, pests, or broken irrigation pipes. These objects are extremly tiny and often take up less than 0.1\% of the total image \cite{agriculture2021}. This is known as the Small Object Detection \gls{sod} problem.

To handle this, agricultural models use a technique called Tiling, also known as Patch based Inference \cite{agriculture2021}, \cite{sahi2022}. Instead of shrinking the whole image, the image is cut into a grid of smaller sections called tiles then Each tile is processed separately which preserves what is called the Ground Sampling Distance \gls{gsd}. In simple terms, GSD refers to how much real world area a single pixel represents. By keeping the tiles at full resolution, no spatiial detail is lost. One well-known implementation is the Slicing Aided Hyper Inference \gls{sahi} framewor\cite{sahi2022}. Studies using SAHI show that tiling can improve small object detection accuracy by 20 to 30\% compared to simply resizing the full image [3].

However, these agricultural solutions have a limitation which is; they are built for powerful workstations and are typically run offline and were not designed for real time processing on small, low powered devices. 



\subsection*{Growing AI and Evolving Edge Computing}
Artificial Intelligence is no longer confined to large data centres in the cloud rather is increasingly happening closer to where data is being collected. This shift is known as edge computing. By 2025, it is estimated that over 75\% of enterprise data is created and processed outside traditional data centres\cite{forbes2025}.
This change is being driven by the need for speed. Applications like autonomous surveillance and industrial monitoring cannot afford the delay of sending data to a remote server and waiting for a response which can be solved by processing locally\cite{assured2025}.

Edge computing works through a network of small, devices ranging from simple IoT components to more sophiscated hardware like NVIDIA Jetson where each device performs some or all of the AI processing locally, without relying on the internet and cloud.

Modern AI models are growing more large and complex and single edge device often does not have enough memory along with processing power to run them. This has led to a new dimension namely Distributed Edge Inference \cite{distredge2025}  where Insteead of running the entire model on one single device, the workload is split across several interconnected devices where each device is responcible for smaller portion of larger workload.

\section{Motivation}
<<<<<<< HEAD
Modern cameras are getting better and better as 4K and 8K sensors are now common in security and industrial monitoring. But AI models have not kept up. Most high performance models still expect small, fixed size inputs. This gap is the core motivation for this research.
Currently, there are two common ways to deal with this mismatch but neither works well.
The first option is to shrink the image size. This is rapid and cheap, but it removes fine details. Small or distant subjects become impossible to detect reliably.
The second option is to use more powerful hardware. \gsls{gpu} can handle large images, but they consume a lot of power. This makes them unsuitable for remote locations or battery powered devices.
So some middle ground is needed. his research aims to build  an architecture that keeps the full detail of the original image, while distributing the processing load across multiple small, affordable devices.


\section{Problem Definition}
This research addresses a core problem in high resolution computer vision. As cameras get better and better, the gap between how much data they capture and speed of processing keeps growing. This is sometimes called the Information Computation Paradox. It breaks down into three specific problems.
\textbf{The Scaling Bottleneck and Receptive Field Constraints:}Most \gls{cnn} models like ResNet or \gls{yolo} expect a fixed input size, typically 416×416 or 640×640 pixels \cite{luo2016understanding} . When a 4K video frame is forced into that size, the image must be heavily resized. This resizing acts like a filter by smoothing out fine details and removes the sharp edges that define small objects. The result is that the region a \gls{cnn} focuses on becomes larger than the object it is trying to find. The network fails to detect it \cite{luo2016understanding}.


\textbf{Hardware and Memory Constraints:}Edge devices are small and have limited memory so running a large AI model directly on one while also holding a high  resolution image in memory is often not possible. The device either crashes with an Out of Memory (OOM) error or slows down due to overheating. There are currently no widely adopted frameworks that allow a single inference task to be split spatially across a group of low power devices \cite{hadidi2019distributed}.
=======
Security and industrial monitoring now frequently use 4k and 8k sensors.As Cameras are getting better and better. But AI models haven’t kept pace. Most high performance models still expect small, fixed sized inputs which is the motivation for this work.
There are two common ways to deal with this mismatch and neither works well.
The first option is to reduce the image size.  It is inexpensive and fast, but loses fine detail. Small or distant objects cannot be reliably detected with this method.
The other option is to use more powerful hardware. The \gls{gpu} can process very big images but they use a lot of compute power. This makes them not suitable for battery operated devices or at remote locations.
So we need some kind of middle ground that address the iisues mentioned above. The goal of his research is to develop an architecture that keeps the detail of the original image, but distributes the processing load among many small and cheap devices.


\section{Problem Definition}

The research works on the fundamentan problem in high res image processing. As camera are getting better and better the amount of pixels they capture and rate of processing is getting larger and larger. Which is popularly known as information computation paradox. It further breaks down into 3 particular problem : 

\textbf{The Scaling Bottleneck and Receptive Field Constraints:} 

Most \gls{cnn} models such as ResNet or \gls{yolo} are designed to take a fixed input size such as 416×416 or 640×640 pixels \cite{luo2016understanding} . That size is going to require massive resizing of the image to fit a 4K video frame. This resizing is essentially a filter that blurs fine detail and removes the hard edges that make small objects pop. As a consequence, the region a \gls{cnn} focuses on becomes larger than the object it is trying to find. The network cannot recognize it \cite{luo2016understanding}.


\textbf{Hardware and Memory Constraints:} Edge devices are small with limited memory, so it is often not possible to run a large AI model directly on one and hold a high resolution image in memory. The device hangs with an Out of Memory (OOM) error or starts lagging due to heat. Currently, there are no common frameworks available that would allow a single inference task to be partitioned in space among a cluster of low power devices \cite{hadidi2019distributed}.
>>>>>>> a0f91c8 (humanized section 1 2 3)

\textbf{The Aggregation Challenge:} Dividing the image into smaller times create a problem near thboundries of each tile. If a person is split among multiple tiles then no simgle device will have complete information to deduce if the object is human or not. The major challange comes to splitting the image and designing the arrgegator algorithm in such a way that it can still make decision based on partial decision from each node without using a much powerful hardware.  \cite{sahi2022_overlap}

\newpage
\section{Objectives}
\begin{itemize}
    \item To build a smaller light weight model that can process a heigher resolution image by dividing it into smaller tiles.
    \item To detect the presence of human by utilizing the above model on a distributed archoitecture of edge devices.
\end{itemize}

\section{Scope and Applications}
This project address the two different dimensions, distributed computing and computer vision. Here Human Presence Detection is taken as a main test case but the model and the system architecture itself is not limited to this particular task only. By making the high resolution image processing more affordable and less complex to precess on edge devices this project targets to solve problems across more than one fields. 


\begin{itemize} 
<<<<<<< HEAD
\item \textbf{High-Precision Industrial Safety:} Large factory floors need constant monitoring. Small but dangerous events, such as a hand getting too close to moving machinery, must be caught early. This system can detect such events at high resolution without the need for expensive industrial computers.
\item \textbf{Privacy-Preserving Smart Homes:}Because each tile is processed locally and only the final results are shared, raw images never leave the device. This makes the system a good fit for home presence sensing and elderly care, where privacy is a major priority.
\item \textbf{Intelligent Energy Management of \gls{hvac} Systems:} Large offices and commercial buildings need accurate occupancy data to manage their Heating, Ventilation, and Air Conditioning \gls{hvac} systems efficiently. Traditional Passive Infrared \gls{pir} sensors struggle to detect people who are sitting still to solve this a distributed CNN based system can detect occupants regardless of whether they are moving or stationary. This information can be fed directly into a Building Management System (BMS), allowing the HVAC system to adjust airflow and temperature room by room in real time.

=======
\item \textbf{High-Precision Industrial Safety:} Large nindustires often need survilance for various issues like hazzards that impact machineries and human life. For example detecting and awaring the worker if hand goes in dangerous zone of moving machine.  


\item \textbf{Privacy-Preserving Smart Homes:} As each image is processing in distributed architecture without needing any connection and datatransfer over vloud this can be used in smart homes without worrying about the data privacy risks. 

\item \textbf{Intelligent Energy Management of \gls{hvac} Systems:} Large commercial buildings use Heating, Ventilation, and Air Conditioning \gls{hvac} systems which are power hungry and require effective managemant. Inorder to effectively manage the \gls{hvac} systems we need the occunpancy data ie. If a human is present ina aroom or not. Traditionally \gls{pir} sensors are being used but they struggle to detect stationary humans. But a image processing system can easily detect presence of humans regardless of motional or not. This information can be used to manage \gls{hvac} systems. 
>>>>>>> a0f91c8 (humanized section 1 2 3)
\end{itemize}
This becomes effective for static or slow moving objects. The system works best in scenes where changes happen slowly enough that the processing delay does not cause missed detections.
